% ===== tr-03 块3：第 3 章 pp.29--31（原文页 29--31）=====
% [原文 p.29]
\begin{enumerate}[label=(\arabic*)]
\item 我们从带丰沛极化的 Calabi--Yau 流形环面退化 $\mathcal X\to D$ 出发。
\item 由这些数据如上所述构造 $(B,\mathcal P,\varphi)$。
\item 作离散 Legendre 变换，得到 $(\check B,\check{\mathcal P},\check\varphi)$。
\item 试着构造 Calabi--Yau 流形的极化退化 $\check{\mathcal X}\to D$，使其对偶相交复形为 $(\check B,\check{\mathcal P},\check\varphi)$。
% [Tutor Note: 第 (4) 步里的 $\check{\mathcal X}$ 是花体 X 带 check，即第 (1) 步中 $\mathcal X$ 的对偶侧对象，不要与前文总空间 $\mathcal X$ 混同；它的分量（如 $\mathcal X_0$）也带 check]
\end{enumerate}

\begin{example}{3.7.10}
离散 Legendre 变换使我们能重现 Batyrev 对偶。回到第 4 节的构造与例 3.7.3：选取 $\tilde\Sigma'$ 上一个相应于线丛 $\mathcal L$ 的严格凸分段线性函数，会诱导 $\mathcal X'_\Delta$ 的一个极化。这便给出 $(B,\mathcal P)$ 上的一个严格凸多值分段线性函数 $\varphi_{\mathcal L}$，从而给出一个离散 Legendre 变换 $(\check B,\check{\mathcal P},\check\varphi_{\mathcal L})$。在 [19] 中我证明了：它是与 $\tilde{\check\Sigma}$ 的某个加细 $\tilde{\check\Sigma}'$ 相对应的对偶相交复形，该加细由在第 4 节的构造中交换 $\nabla$ 与 $\Delta$ 的角色而得到。
% [Tutor Note: 本句出现两套记号：开头的 $\tilde\Sigma'$ 是 $\Delta$ 侧（p.24）的扇；末尾的 $\tilde{\check\Sigma}'$ 与 $\tilde{\check\Sigma}$ 是交换 $\nabla$ 与 $\Delta$ 后 $\nabla$ 侧的扇（tilde 加 check 的双重记号）。二者是不同对象，并非同一对象前后不一致；按原文分别照排]
% [Tutor Note: 这里 $\tilde{\check\Sigma}$ 是双重复合记号：紧贴 Σ 的是 check（ˇ），其上方还有 tilde（~），不要把上一层的 tilde 漏掉、写成 $\Delta$ 侧的 $\tilde\Sigma$；1200 dpi 裁图（p29 该行）可清晰分辨两层重音]
作为练习，读者可以自行验证下述事实。若取 $\check h=\check\varphi$，并另外定义
\[
h=\varphi:N_\R\to\R
\]
使它在 $\Sigma$（即 $\Delta$ 的正规扇）的每个一维锥的本原生成元上取值 1，则由第 4 节可在 $B=\partial\nabla$ 上得到带奇点的仿射结构；完全对称地，用 $h$ 也可在 $\check B=\partial\Delta$ 上得到这样的结构。这些流形分别带有由 $\nabla$ 与 $\Delta$ 的所有真面构成的多面体分解 $\mathcal P$ 与 $\check{\mathcal P}$。$\PP_\Delta$ 与 $\PP_\nabla$ 上的反典范极化分别诱导 $B$ 与 $\check B$ 上的多值分段线性函数 $\psi$、$\check\psi$。然后证明 $(B,\mathcal P,\psi)$ 与 $(\check B,\check{\mathcal P},\check\psi)$ 互为对方的离散 Legendre 变换。
% [Tutor Note: 原文此处作 completely symetrically，系 symmetrically 的拼写笔误（少一个 m）；本译本按"完全对称地"译出]
\end{example}

于是 Batyrev（以及 Batyrev--Borisov）对偶是本构造的一个特例。

这个镜像对称算法中唯一缺失的是最后一步：

\begin{question}{3.7.11}
（重构问题，版本 II）给定 $(B,\mathcal P,\varphi)$，能否构造一个极化环面退化 $\mathcal X\to D$，使其对偶相交复形为 $(B,\mathcal P,\varphi)$？
\end{question}

从数据 $(\check B,\check{\mathcal P},\check\varphi)$ 重构一个退化的中心纤维 $\mathcal X_0$ 是相当显然的，我们将在第 8 节中明确看到这一点。人们或许会天真地希望：这个可约簇具有良好形变理论，从而能被光滑化。然而一般说来，它的形变理论性质很坏。

正如 Kawamata 与 Namikawa 在 [36] 中对正规交叉情形的最初观察：需要在 $\mathcal X_0$ 上附加某种额外结构，它才具有良好的形变理论。这个结构就是对数结构；引入对数结构使我们能直接在退化纤维本身上研究镜像对称的许多方面。

我们将在下一节做这件事；但首先，让我来讨论这个镜像对称构造可能有多一般的问题：

\begin{conjecture}{3.7.12}
若 $f:\mathcal X\to D$ 是大复结构极限退化，则 $f$ 双有理等价于一个环面退化 $f':\mathcal X'\to D$。
\end{conjecture}

大复结构极限这一条件，按 Morrison 在 [47] 中的定义，比极大幂零更强。我为何会想象这样的事情成立呢？

% [原文 p.30]
好，自由地幻想一下：我们应当期望，在纤维 $X_t$ 上选定极化与 Ricci 平坦度量之后，有序列 $(X_t,g_t)$ 收敛到一个带奇点的仿射流形 $B$。然而一般说来，带奇点的仿射流形未必能作为某个环面退化的对偶相交复形出现；首先，它未必有多面体分解。举例来说，即使在二维情形也存在轨形奇点（对应于非半稳定的奇异椭圆纤维），它们不出现在环面退化的对偶相交复形中，却可以作为 K3 曲面上某个特殊拉格朗日纤维化的底。不过 K3 情形下，一般的底确实能作为某个环面退化的对偶相交复形出现。希望在于：大复结构极限这一条件迫使 $B$ 的奇点"足够一般"，从而使人们能在 $B$ 上构造一个良好的环面多面体分解 $\mathcal P$，并由此构造出一个环面退化。若选择得当，这个环面退化想必会与原退化双有理等价。
% [Tutor Note: 上文 orbifold singularities 译作"轨形奇点"（orbifold 的通行译名之一），指底空间局部为有限群商的工作方式产生的奇点；它作为 SYZ 底空间出现时并不来自环面退化的对偶相交复形]

这一论证当然是相当粗疏的（hand-wavy），但我相信它提供了某种道义上的期望：可能存在一大类退化，我们的方法对之适用。我还要指出，在 K3 曲面的情形可以证明该猜想。
% [Tutor Note: hand-wavy 是口语表达，意为"粗略、不严格、靠手势比划式的"；此处按其语气译作"粗疏"。该段其余部分（"道义上的期望"= moral expectation）也是作者口语化的措辞，非术语]

我们现在来到 [22] 所铺陈纲领的技术核心。这个纲领的某些方面相当技术化，因此这里的目标是尽可能简单地解释 [22] 的要点。
% [Tutor Note: 原文作 explain the highlights of [22] as simply possible，系 as simply as possible 之省字笔误；本译本按"尽可能简单地"译出]

\section{对数结构}

我们首先介绍 Fontaine--Illusie 与 Kato 的对数结构（[32]、[35]）。

\begin{definition}{3.8.1}
概形（或解析空间）$X$ 上的一个对数结构，是（乘法、交换的）幺半群层之间的一个（含幺）同态
\[
\alpha_X:\mathcal M_X\to\mathcal O_X,
\]
它诱导出同构 $\alpha_X^{-1}(\mathcal O_X^\times)\to\mathcal O_X^\times$。三元组 $(X,\mathcal M_X,\alpha_X)$ 于是称为一个\emph{对数空间}。我们常把整个包裹写成 $X^\dagger$。
% [Tutor Note: 定义中的 $\alpha_X^{-1}(\mathcal O_X^\times)$ 指层的前像（preimage），不是 $\alpha_X$ 的倒数；同构的方向是 $\alpha_X^{-1}(\mathcal O_X^\times)\to\mathcal O_X^\times$，不要写成反向]
\end{definition}

对数空间之间的态射 $F:X^\dagger\to Y^\dagger$，由一个底空间之间的态射 $\underline F:X\to Y$，连同与结构同态交换的同态 $F^\#:\underline F^{-1}(\mathcal M_Y)\to\mathcal M_X$ 组成：
\[
\alpha_X\circ F^\#=\underline F^*\circ\alpha_Y.
\]
% [Tutor Note: 上式中 $\underline F$（带下划线）表示底空间之间的态射，以区别于对数空间之间的态射 $F$；$\underline F^*$ 是它对函数的拉回。原文未解释下划线的含义，本译本亦不额外加注于正文]

关键的几类例子如下。

\begin{example}{3.8.2}
\begin{enumerate}[label=(\arabic*)]
\item 设 $X$ 是概形，$Y\subseteq X$ 是余维 1 的闭子集。以 $j:X\setminus Y\to X$ 记此包含。则包含
\[
\alpha_X:\mathcal M_X=j_*(\mathcal O_{X\setminus Y}^\times)\cap\mathcal O_X\to\mathcal O_X
\]
（即零点含于 $Y$ 中的正则函数层）是 $X$ 上的一个对数结构。它称为 $X$ 上的\emph{除子对数结构}。
\item 一个\emph{预对数结构}，即幺半群层之间的任意同态 $\varphi:\mathcal P\to\mathcal O_X$，如下定义一个相伴的对数结构 $\mathcal M_X$：
\[
\mathcal M_X=(\mathcal P\oplus\mathcal O_X^\times)/\{(p,\varphi(p)^{-1})\mid p\in\varphi^{-1}(\mathcal O_X^\times)\},
\]
且 $\alpha_X(p,h)=h\cdot\varphi(p)$。
% [原文 p.31]
\item 若 $f:X\to Y$ 是概形态射，且 $\alpha_Y:\mathcal M_Y\to\mathcal O_Y$ 是 $Y$ 上的对数结构，则由预对数结构 $f^{-1}(\mathcal M_Y)\to\mathcal O_X$ 定义出 $X$ 上一个相伴的对数结构，即\emph{拉回对数结构}。
\item 在 (1) 中，我们可以用 (3) 把 $X$ 上的对数结构拉回到 $Y$。于是特别地，若 $\mathcal X\to D$ 是环面退化，则包含 $\mathcal X_0\subseteq\mathcal X$ 给出 $\mathcal X$ 上的一个对数结构，并诱导出 $\mathcal X_0$ 上的一个对数结构。类似地，包含 $0\in D$ 给出 $D$ 上的一个对数结构，并诱导出 $0$ 上的一个对数结构。这里 $\mathcal M_0=\C^\times\oplus\N$，其中 $\N$ 是自然数（非负整数）的（加法）幺半群，而
\[
\alpha_0(h,n)=\begin{cases} h & n=0,\\ 0 & n\ne0.\end{cases}
\]
$0^\dagger$ 通常称为\emph{标准对数点}。于是我们有对数态射 $\mathcal X^\dagger\to D^\dagger$ 与 $\mathcal X_0^\dagger\to 0^\dagger$。
\item 设 $\sigma\subseteq M_\R=\R^n$ 是凸有理多面体锥，$\sigma^\vee\subseteq N_\R$ 是其对偶锥，置 $P=\sigma^\vee\cap N$：这是一个幺半群。由 $\sigma$ 定义的仿射环面簇可写成 $X=\operatorname{Spec}\C[P]$。于是由幺半群同态
\[
P\to\C[P],\qquad p\mapsto z^p
\]
诱导出一个预对数结构，从而在 $X$ 上有一个相伴的对数结构。这事实上与由 $\partial X\subseteq X$ 诱导的对数结构相同，其中 $\partial X$ 是 $X$ 的环面边界，即 $X$ 的所有环面除子之并。

若 $p\in P$，则单项式 $z^p$ 定义一个映射 $f:X\to\operatorname{Spec}\C[\N]$（$=\operatorname{Spec}\C[t]$），它是一个对数态射，其中 $\operatorname{Spec}\C[\N]$ 上的对数结构由 $\N\to\C[\N]$ 类似地诱导。纤维 $X_0=\operatorname{Spec}\C[P]/(z^p)$ 是 $X$ 的子概形，其上有诱导的对数结构，并如 (4) 那样有映射 $X_0^\dagger\to0^\dagger$。$f$ 是\emph{对数光滑}态射的一个例子。本质上所有对数光滑态射都平展局部地具有这种形式（若把 $\N$ 换成更一般的幺半群）。细节见 [34]。
\end{enumerate}

定义 3.7.1 的条件 (4) 事实上蕴涵：局部地、在远离 $Z$ 处，$\mathcal X^\dagger$ 与 $\mathcal X_0^\dagger$ 具有上述形式。因此我们应把 $\mathcal X^\dagger\to D^\dagger$ 看作在远离 $Z$ 处对数光滑的；而从对数的观点看，$\mathcal X_0^\dagger$ 可以处理得很像一个远离 $X$ 的非奇异概形。
% [Tutor Note: 末句原文作 away from X，但两句话之前刚说 away from Z，且 $X$ 在此处并非被排除的对象，疑为 $Z$ 之笔误；正文照原文译作"远离 $X$"（550 dpi 裁图确认印作 X）]
我们将在下面谈到微分时明确看到这一点。\proofend
\end{example}

在对数概形 $X^\dagger$ 上总存在正合列
\[
1\to\mathcal O_X^\times\xrightarrow{\alpha^{-1}}\mathcal M_X\to\overline{\mathcal M}_X\to0,
\]
其中我们把幺半群商层 $\overline{\mathcal M}_X$ 用加法记号书写。我们称 $\overline{\mathcal M}_X$ 为该对数结构的\emph{幽灵层}。我喜欢把 $\overline{\mathcal M}_X$ 看作指定了联系于该对数结构的组合信息。例如，若 $X^\dagger$ 由 Cartier 除子 $Y\subseteq X$ 诱导（$X$ 正规），则 $\overline{\mathcal M}_{X,x}$ 在 $x\in X$ 处的茎，是 $x$ 的某个邻域上支在 $Y$ 中的有效 Cartier 除子构成的幺半群。

\begin{exercise}{3.8.3}
证明：在例 3.8.2 的 (5) 中，若 $\dim\sigma=n$ 且 $x$ 是 $X$ 唯一的零维环面轨道，则 $\overline{\mathcal M}_{X,x}=P$。更一般地，
\[
\overline{\mathcal M}_{X,x}=\frac{\tau^\vee\cap N}{\tau^\perp\cap N}=\operatorname{Hom}_{monoid}(\tau\cap M,\N),
\]
其中 $x\in X$ 位于 $\sigma$ 的面 $\tau$ 所对应的环面轨道中。特别地，$\tau$ 可由 $\operatorname{Hom}_{monoid}(\overline{\mathcal M}_{X,x},\R_{\ge0}^+)$ 恢复出来，其中 $\R_{\ge0}^+$ 是非负实数的加法幺半群。\proofend
\end{exercise}
